% Pista ana-26s-q4a · Anàlisi
% Procedència: PAU Catalunya 2026 · Convocatòria de setembre · Sèrie 2 · Exercici 4 — Opció A
\documentclass[11pt,a4paper]{article}
\usepackage{pista}
\pistainfo{Catalunya 2026}{Convocatòria de setembre}{Sèrie 2}

\begin{document}

\section*{\textcolor{blauPista}{Exercici 4 --- Opció A}}

\noindent L'Aina vol construir una capsa de cartró amb forma de prisma rectangular amb tapa, per guardar-hi la tauleta i el carregador.

\noindent L'Aina construirà la capsa a partir d'un cartró rectangular de $80\,\text{cm} \times 50\,\text{cm}$. A la figura següent es mostra el desplegament de la capsa; les zones ratllades són les parts que l'Aina retallarà i retirarà per tal de poder-la muntar:

\subsection*{4.1. \normalfont Comproveu que el volum de la capsa resultant és $V(x) = 2x^{3} - 130x^{2} + 2\,000x$, on $x$ és la longitud del costat del quadrat que cal retallar dels cantons per a poder-la muntar. \hfill\textnormal{[1~punt]}}

\pista{Mira el desplegament al llarg dels $80\,\text{cm}$: d'esquerra a dreta hi ha una franja d'amplada $x$ (un lateral), la base, una altra franja d'amplada $x$ (l'altre lateral) i la tapa, que ha de ser igual que la base. Compte: no és la capsa oberta de sempre (on la llargada seria el costat del cartró menys $2x$), perquè el que queda dels $80\,\text{cm}$ s'ha de repartir entre la base i la tapa. En vertical, fixa't en quantes franges d'amplada $x$ s'hi retiren. Escriu la llargada, l'amplada i l'alçada en funció de $x$, multiplica-les i desenvolupa el producte fins a arribar a l'expressió de l'enunciat.}

\subsection*{4.2. \normalfont Calculeu les dimensions de la capsa per tal que el seu volum sigui màxim. Quin és aquest volum màxim? \hfill\textnormal{[1,5~punts]}}

\pista{Deriva $V(x)$ i iguala la derivada a zero: obtindràs una equació de segon grau (pots dividir-la abans per $2$) amb dues solucions. No te les quedis totes dues: pensa quins valors de $x$ tenen sentit, és a dir, per a quins les tres dimensions de l'apartat 4.1 són positives; això et permet descartar-ne una. Comprova que l'altra és un màxim, per exemple amb el signe de $V''(x)$. Finalment, rellegeix la pregunta: et demanen les tres dimensions de la capsa (no només el valor de $x$) i el volum màxim, amb les unitats ($\text{cm}$ i $\text{cm}^{3}$).}

\end{document}
